\chapter{Probability Generating Functions}\label{ch:b2:genfun}

The power series of \cref{ch:b2:powerseries} return with a
probabilistic mission: to an $\N$-valued random variable we attach
the power series with coefficients $\P(X = n)$. This
\emph{generating function} converts sums of independent variables
into products, moments into derivatives at $1$, and hard
combinatorial identities into one-line multiplications. The
chapter closes the book with two showpieces: the Poisson
approximation of rare events, and the extinction criterion for
branching processes --- a genuinely infinite probabilistic
computation solved entirely by the geometry of a convex curve.

\section{Definition and basic properties}

\begin{definition}[Probability generating function]\label{def:b2:genfun:pgf}
Let $X$ be an $\N$-valued random variable, $p_n = \P(X = n)$. The
\emph{(probability) generating function} of $X$ is the sum of the
power series
\[
G_X(t) = \E\bigl(t^X\bigr) = \sum_{n=0}^{\infty} p_n\,t^n .
\]
\end{definition}

\begin{example}[First reflexes]\label{ex:b2:genfun:firstreflexes}
A constant variable $X = c$ has $G_X(t) = t^c$; a shift obeys
$G_{X+c}(t) = t^c\,G_X(t)$; and evaluating at special points
reads off information without any expansion: $G_X(0) = \P(X
= 0)$, $G_X(1) = 1$, and $G_X(-1) = \P(X\text{ even}) -
\P(X\text{ odd})$, the parity balance exploited in
\cref{exo:b2:genfun:10}. These one-liners are used silently
everywhere below --- and the evaluation $G_X(0)$ is exactly
how extinction probabilities will be extracted from iterated
generating functions at the end of the chapter.
\end{example}

\begin{proposition}[Radius and first properties]\label{prop:b2:genfun:radius}
The series defining $G_X$ has radius of convergence $\geq 1$;
$G_X$ is defined and continuous on $\intcc{-1}{1}$,
$\mathcal{C}^\infty$ on $\intoo{-1}{1}$, with $G_X(1) = 1$ and
$\abs{G_X(t)} \leq 1$ there. Moreover $G_X$ determines the law of
$X$:
\[
p_n = \frac{G_X^{(n)}(0)}{n!} .
\]
\end{proposition}

\begin{proof}
Since $\sum p_n = 1$ converges, the terms $p_n\,1^n$ are bounded,
so the radius is $\geq 1$ (Abel's lemma,
\cref{ch:b2:powerseries}); at $t = \pm1$ the series converges
absolutely ($\sum p_n = 1$ dominates); better, on the whole
interval $\intcc{-1}1$,
\[
\sup_{\abs t\leq1}\,\abs{p_nt^n} = p_n
\quad\text{with}\quad \sum_np_n < \infty :
\]
the series converges \emph{normally} on $\intcc{-1}1$, so
its sum is continuous there
(\cref{thm:b2:funcseq:weierstrass,thm:b2:funcseq:continuity}). Smoothness inside and the coefficient
formula are the general theory of power series; the coefficients
being recoverable, two variables with the same generating function
have the same law.
\end{proof}

\begin{example}[The classical laws]\label{ex:b2:genfun:classical}
\leavevmode
\begin{itemize}
\item Bernoulli $\mathcal{B}(p)$: $G(t) = 1 - p + pt$.
\item Binomial $\mathcal{B}(n, p)$: $G(t) = \sum_k \binom nk
(pt)^k(1-p)^{n-k} = (1 - p + pt)^n$ (binomial theorem).
\item Geometric $\mathcal{G}(p)$: $G(t) =
\sum_{k\geq1}(1-p)^{k-1}p\,t^k = \dfrac{pt}{1 - (1-p)t}$ (radius
$\frac{1}{1-p} > 1$).
\item Poisson $\mathcal{P}(\lambda)$: $G(t) =
\sum_k e^{-\lambda}\frac{(\lambda t)^k}{k!} = e^{\lambda(t - 1)}$
(radius $\infty$).
\end{itemize}
\end{example}

\begin{example}[Integrating the generating
function]\label{ex:b2:genfun:integraltrick}
Derivatives of $G_X$ at $1$ give positive moments; the
\emph{integral} gives a negative one. From $\int_0^1t^k\dd t
= \frac1{k+1}$ and term-by-term integration (normal
convergence on $\intcc01$):
\[
\int_0^1G_X(t)\,\dd t = \sum_{k\geq0}\frac{\P(X =
k)}{k+1} = \E\Bigl(\frac1{1+X}\Bigr).
\]
For $X \sim \mathcal P(\lambda)$:
\[
\E\Bigl(\frac1{1+X}\Bigr) =
\int_0^1\eu^{\lambda(t-1)}\,\dd t = \frac{1 -
\eu^{-\lambda}}{\lambda},
\]
recovering in one line the series computation of
\cref{ex:b2:randomvar:transferex}. The generating function
is a two-way instrument: differentiate at $1$ for the
moments $\E(X)$, $\E(X(X-1))$, integrate over $\intcc01$
for $\E\bigl(\frac1{1+X}\bigr)$ --- one analytic object,
queried in whichever direction the problem needs.
\end{example}

\begin{example}[A law with radius exactly
one]\label{ex:b2:genfun:heavytail}
Let $\P(X = k) = \dfrac{6}{\pi^2k^2}$ for $k \geq 1$ --- a
probability law by Basel's identity
(\cref{ex:b2:fourier:basel}). Its generating function
$G(t) = \frac6{\pi^2}\sum_{k\geq1}\frac{t^k}{k^2}$ has radius
of convergence exactly $1$: the general bound ``radius $\geq
1$'' of \cref{prop:b2:genfun:radius} cannot be improved. And
the mean is
\[
\sum_{k\geq1}k\,\P(X = k) =
\frac6{\pi^2}\sum_{k\geq1}\frac1k = \infty :
\]
$G$ is continuous on $\intcc{-1}1$, smooth inside, but its
derivative blows up at $1^-$ --- the graph arrives at the
point $(1, 1)$ with a vertical tangent. Heavy tails are
visible \emph{geometrically} on the generating function, at
the single point $t = 1$; the moments theorem below makes
this correspondence exact.
\end{example}

\begin{theorem}[Moments from the generating
function]\label{thm:b2:genfun:moments}
$X$ has an expectation if and only if $G_X$ is differentiable at
$1^-$ (left derivative, finite), and then $\E(X) = G_X'(1)$.
Similarly $X$ has a second moment iff $G_X$ is twice
differentiable at $1^-$, and then
\[
\E\bigl(X(X - 1)\bigr) = G_X''(1),
\qquad
V(X) = G_X''(1) + G_X'(1) - G_X'(1)^2 .
\]
\end{theorem}

\begin{proof}
For $t \in \intoo{0}{1}$, term-by-term differentiation inside the
disk gives $G_X'(t) = \sum_{n\geq1} np_n t^{n-1}$, a series with
nonnegative coefficients: $t \mapsto G_X'(t)$ is nondecreasing on
$\intoo{0}{1}$, and by monotone convergence of partial sums (or
Abel's theorem for nonnegative coefficients,
\cref{ch:b2:powerseries}),
\[
\lim_{t \to 1^-} G_X'(t)
= \sum_{n\geq1} n\,p_n \in \intcc{0}{+\infty} ,
\]
each side finite exactly when the other is. When finite, the mean
value theorem squeezes the difference quotients $\frac{G_X(1) -
G_X(t)}{1 - t}$ between values of $G_X'$, so $G_X$ is
differentiable at $1^-$ with $G_X'(1) = \sum np_n = \E(X)$ (by
transfer). The second-order statement repeats the argument one level
up: $G''_X(t) = \sum_{n\geq2}n(n-1)p_nt^{n-2}$ is
nondecreasing on $\intoo01$ with monotone limit
$\sum_nn(n-1)p_n = \E(X(X-1))$, finite exactly when $X$ has
a second moment. The variance formula then follows from
K\"onig--Huygens:
\[
V(X) = \E(X^2) - \E(X)^2 = \E\bigl(X(X-1)\bigr) + \E(X) -
\E(X)^2 = G''_X(1) + G'_X(1) - G'_X(1)^2 .
\]
\end{proof}

\begin{example}
Poisson: $G'(t) = \lambda e^{\lambda(t-1)}$, so $\E(X) = \lambda$;
$G''(1) = \lambda^2$, so $V(X) = \lambda^2 + \lambda - \lambda^2 =
\lambda$ --- the computations of \cref{ch:b2:randomvar} in one
line each.
\end{example}

\begin{example}[The mode of a Poisson law]
\label{ex:b2:genfun:poissonmode}
Where is $\P(X = k)$ largest for $X \sim \mathcal
P(\lambda)$? Consecutive weights compare through the ratio
\[
\frac{\P(X = k+1)}{\P(X = k)} = \frac{\lambda}{k + 1} ,
\]
which exceeds $1$ while $k < \lambda - 1$ and drops below
$1$ once $k > \lambda - 1$: the weights rise then fall, with
mode $\floor\lambda$ (and a tie between $\lambda - 1$ and
$\lambda$ when $\lambda$ is an integer: for $\lambda = 3$,
$\P(X = 2) = \P(X = 3) = \frac92\eu^{-3} \approx 0.224$).
Ratio tests on the coefficients are often the fastest route
to qualitative facts about a discrete law --- no generating
function needed, but the coefficients \emph{are} the
generating function, read term by term.
\end{example}

\section{Sums of independent variables}

\begin{theorem}[Multiplicativity]\label{thm:b2:genfun:product}
If $X$ and $Y$ are independent $\N$-valued random variables, then
\[
G_{X + Y}(t) = G_X(t)\,G_Y(t)
\qquad (\abs t \leq 1),
\]
and by induction $G_{X_1 + \dots + X_n} = \prod_i G_{X_i}$ for
independent $X_1, \dots, X_n$.
\end{theorem}

\begin{proof}
Two proofs, both instructive. \emph{Via expectations:} $t^X$ and
$t^Y$ are independent bounded variables, so
(\cref{thm:b2:randomvar:product})
\[
G_{X+Y}(t) = \E\bigl(t^{X+Y}\bigr)
= \E\bigl(t^X t^Y\bigr)
= \E\bigl(t^X\bigr)\E\bigl(t^Y\bigr) .
\]
\emph{Via Cauchy products:} the law of $X + Y$ is the convolution
$\P(X + Y = n) = \sum_{k=0}^n \P(X = k)\P(Y = n - k)$, and the
Cauchy product theorem for absolutely convergent series
(\cref{ch:b2:series}) multiplies the two power series exactly
along this convolution.
\end{proof}

\begin{example}[Stability of the classical laws]\label{ex:b2:genfun:stability}
Independent binomials with the same $p$ add: $(1 - p + pt)^m(1 -
p + pt)^n = (1 - p + pt)^{m+n}$, so $\mathcal{B}(m, p) +
\mathcal{B}(n, p) = \mathcal{B}(m + n, p)$ --- in particular a sum
of $n$ independent Bernoulli variables is binomial, re-proving the
law of the number of successes. Independent Poissons add:
$e^{\lambda(t-1)}e^{\mu(t-1)} = e^{(\lambda + \mu)(t-1)}$, so
$\mathcal{P}(\lambda) + \mathcal{P}(\mu) = \mathcal{P}(\lambda +
\mu)$ --- the convolution computation of
\cref{exo:b2:randomvar:2}, now without computation.
\end{example}

\begin{example}[Two dice, one polynomial
squared]\label{ex:b2:genfun:twodice}
For one fair die, $G(t) = \frac{t + t^2 + \dots + t^6}{6}$;
for the sum of two,
\[
G(t)^2 = \frac{1}{36}\bigl(t^2 + 2t^3 + 3t^4 + 4t^5 + 5t^6 +
6t^7 + 5t^8 + 4t^9 + 3t^{10} + 2t^{11} + t^{12}\bigr) :
\]
the triangular law of dice sums ($7$ is the mode, with
probability $\frac6{36} = \frac16$), read off a polynomial
square that one multiplies out once in a lifetime. The
convolution formula would have required eleven separate
counting arguments; the generating function does them all
simultaneously, because multiplying polynomials \emph{is}
convolving coefficients. This mechanical translation ---
laws to coefficients, sums to products --- is the whole
business model of the chapter, and
\cref{exo:b2:genfun:11} pushes it to the surprising
Sicherman dice.
\end{example}

\begin{example}[Three dice and a coefficient
extraction]\label{ex:b2:genfun:threedice}
For the sum $S$ of three fair dice, $\P(S = 10)$ is the
coefficient of $t^{10}$ in $\bigl(\frac{t + \dots +
t^6}6\bigr)^3$. Factor and expand with the binomial and
geometric series:
\[
\Bigl(\frac{t(1 - t^6)}{6(1 - t)}\Bigr)^{\!3}
= \frac{t^3}{216}\,\bigl(1 - 3t^6 + 3t^{12} -
t^{18}\bigr)\sum_{j\geq0}\binom{j+2}2t^j .
\]
The coefficient of $t^{10}$ requires $t^7$ from the product:
$j = 7$ with the term $1$, and $j = 1$ with the term
$-3t^6$:
\[
\P(S = 10) = \frac{1}{216}\Bigl(\binom92 -
3\binom32\Bigr) = \frac{36 - 9}{216} = \frac{27}{216} =
\frac18 .
\]
Direct enumeration of the $27$ triples is error-prone; the
algebra is mechanical and scales to any number of dice ---
the inclusion--exclusion visible in $(1 - t^6)^3$ is doing
the case analysis automatically.
\end{example}

\begin{example}[Reading a law off its generating
function]\label{ex:b2:genfun:identify}
Which law has $G(t) = \dfrac1{2 - t}$? Expand into a power
series:
\[
\frac{1}{2 - t} = \frac12\cdot\frac1{1 - t/2}
= \sum_{k\geq0}\frac{t^k}{2^{k+1}} :
\]
nonnegative coefficients summing to $G(1) = 1$, so this is a
genuine law, $\P(X = k) = 2^{-(k+1)}$ on $\N$ --- a
geometric law starting at $0$. By uniqueness
(\cref{prop:b2:genfun:radius}), no other law shares this
$G$. Recognizing laws from their generating functions is a
skill worth drilling: it is how the critical branching
iterate $G_n(t) = \frac{n - (n-1)t}{n+1 - nt}$ of the
weekend problem is unmasked as a geometric law conditioned
on survival.
\end{example}

\begin{remark}
Stability goes one way only: sums of independent Poissons are
Poisson, but differences are not --- $X - Y$ takes negative
values, so it has no generating function at all, and its law
(the Skellam distribution) lies outside this chapter's
toolkit. Likewise $\mathcal B(m, p) + \mathcal B(n, p')$ with
$p \neq p'$ is \emph{not} binomial: the product $(1 - p +
pt)^m(1 - p' + p't)^n$ has two distinct root locations, while
every binomial pgf has a single repeated root. Reading
stability off root patterns is a small preview of how much
structure the polynomial encodes.
\end{remark}

\begin{remark}[The roots-of-unity filter]
Evaluating at $-1$ separates even from odd; evaluating at all
$m$-th roots of unity separates every residue class: with
$\omega = \eu^{2\iu\pi/m}$,
\[
\P(X \equiv r \bmod m)
= \frac1m\sum_{j=0}^{m-1}\omega^{-jr}\,G_X(\omega^j),
\]
since averaging $\omega^{j(k-r)}$ over $j$ yields $1$ if $k
\equiv r$ and $0$ otherwise. Sample dividend: for the sum $S$
of two fair dice, each $G(\omega^j) = \frac16\sum_{k=1}^6
\omega^{jk} = -\frac16$ for $j \neq 0$ (the seven seventh
roots of unity sum to zero), so
\[
\P(7 \mid S) = \frac17\Bigl(1 +
6\cdot\frac1{36}\Bigr) = \frac16 ,
\]
confirming the count from \cref{ex:b2:genfun:twodice} ---
and the method scales to questions where direct counting
does not.
\end{remark}

\begin{theorem}[Random sums: Wald's identity for generating
functions]\label{thm:b2:genfun:compound}
Let $(X_k)_{k\geq1}$ be independent $\N$-valued variables with the
same law and generating function $G_X$, and let $N$ be an
$\N$-valued variable independent of the $X_k$, with generating
function $G_N$. Then the random sum $S = X_1 + \dots + X_N$ (with
$S = 0$ when $N = 0$) has generating function
\[
G_S = G_N \circ G_X .
\]
In particular, if $N$ and $X_1$ have expectations, $\E(S) =
\E(N)\,\E(X_1)$.
\end{theorem}

\begin{proof}
Condition on $N$ (total probability,
\cref{thm:b2:proba:bayes}): for $\abs t \leq 1$,
\[
G_S(t) = \sum_{n=0}^\infty \P(N = n)\,
\E\bigl(t^{X_1 + \dots + X_n}\bigr)
= \sum_{n=0}^\infty \P(N = n)\,G_X(t)^n
= G_N\bigl(G_X(t)\bigr),
\]
using multiplicativity for each fixed $n$ and the summability of
the whole double family ($\abs{G_X(t)} \leq 1$). The exchange of
summations is Fubini for summable families
(\cref{ch:b2:series}). Differentiating at $1^-$ by the chain rule
and \cref{thm:b2:genfun:moments}: $\E(S) = G_N'(G_X(1))\,G_X'(1)
= G_N'(1)G_X'(1) = \E(N)\E(X_1)$.
\end{proof}

\begin{example}[Compound Poisson: yearly insurance
losses]\label{ex:b2:genfun:compoundpoisson}
An insurer receives $N \sim \mathcal P(\lambda)$ claims in a
year, each claim costing $X_k$ (integer units, i.i.d., pgf
$G_X$, mean $\mu$, independent of $N$). By
\cref{thm:b2:genfun:compound}, the total loss $S$ has
\[
G_S(t) = \eu^{\lambda(G_X(t) - 1)},
\qquad
\E(S) = \lambda\mu ,
\]
and differentiating twice at $1^-$:
\[
V(S) = \lambda\,G_X''(1) + \lambda^2\mu^2 + \lambda\mu -
(\lambda\mu)^2 = \lambda\,\E(X^2) .
\]
The variance involves the \emph{second} moment of a single
claim, not its variance: a compound Poisson sum feels the
occasional large claim twice --- once through how many, once
through how big. For $\lambda = 10$ claims of geometric law
with mean $2$ ($\E X^2 = 6$): $\E S = 20$, $V(S) = 60$, and
Chebyshev (\cref{ch:b2:randomvar}) already yields usable
solvency margins. This ``randomly stopped sum'' pattern is the
same one that will drive the branching recursion of
\cref{prop:b2:genfun:branching}: composition of generating
functions is the algebra of random populations.
\end{example}

\begin{remark}
The independence of $N$ from the summands is not decorative.
Take $X_k \in \{0, 2\}$ with equal probabilities and let $N =
X_1$ (flagrantly dependent): then $S = X_1 + \dots + X_N$ is
$0$ when $X_1 = 0$, and $2 + X_2$ when $X_1 = 2$, so $\E(S) =
\frac12(2 + 1) = \frac32$, while $\E(N)\E(X_1) = 1\cdot1 =
1$: Wald's identity fails. When the number of terms is
allowed to \emph{react} to the terms themselves, the clean
product structure collapses --- the full theory of such
``stopping'' rules is the martingale chapter of the Year 3
volume.
\end{remark}

\section{Poisson approximation}

\begin{theorem}[Law of rare events]\label{thm:b2:genfun:rareevents}
Let $X_n \sim \mathcal{B}(n, p_n)$ with $n\,p_n \to \lambda > 0$.
Then for every $k \in \N$:
\[
\P(X_n = k)
\xrightarrow[n\to\infty]{}
e^{-\lambda}\frac{\lambda^k}{k!} :
\]
the binomial law of many rare independent events converges to the
Poisson law of parameter $\lambda$.
\end{theorem}

\begin{proof}
Direct computation with $p_n = \frac{\lambda_n}{n}$, $\lambda_n
\to \lambda$:
\[
\P(X_n = k)
= \binom nk p_n^k(1 - p_n)^{n-k}
= \frac{n(n-1)\cdots(n-k+1)}{n^k}\cdot
\frac{\lambda_n^k}{k!}\,
\bigl(1 - \tfrac{\lambda_n}{n}\bigr)^{n-k} .
\]
As $n \to \infty$ with $k$ fixed: the first factor tends to $1$
(product of $k$ factors $\to 1$); $\lambda_n^k \to \lambda^k$;
and $\bigl(1 - \frac{\lambda_n}{n}\bigr)^{n-k} =
\exp\bigl((n-k)\ln(1 - \frac{\lambda_n}{n})\bigr) \to
e^{-\lambda}$ since $(n - k)\ln\bigl(1 - \frac{\lambda_n}{n}\bigr)
\sim -\lambda_n \to -\lambda$
(\cref{ch:b2:comparison}). Alternatively, at the level of
generating functions: $G_{X_n}(t) = \bigl(1 +
\frac{\lambda_n(t-1)}{n}\bigr)^n \to e^{\lambda(t - 1)} =
G_{\mathcal{P}(\lambda)}(t)$ for each fixed $t \in [0, 1]$ ---
convergence of generating functions, which (for $\N$-valued
variables) is equivalent to convergence of each $\P(X_n = k)$;
see \cref{exo:b2:genfun:9}.
\end{proof}

\begin{remark}
This is why Poisson laws model counts of rare events --- typos per
page, radioactive decays per second, accidents per day at a
crossing: each opportunity is nearly negligible, opportunities are
many, and only the mean rate $\lambda$ survives in the limit.
\end{remark}

\begin{example}[Watching the Poisson limit
converge]\label{ex:b2:genfun:poissonnumerics}
Fix $\lambda = 2$ and let $X_n \sim \mathcal B(n, 2/n)$. The
no-event probability is exactly $\P(X_n = 0) = (1 - 2/n)^n$:
\[
n = 10:\ 0.107, \qquad
n = 20:\ 0.122, \qquad
n = 50:\ 0.130, \qquad
n = 100:\ 0.133,
\]
against the limit $\eu^{-2} \approx 0.135$. The convergence
is monotone and of speed $O(1/n)$ --- expanding, $(1 -
2/n)^n = \eu^{-2}\bigl(1 - \tfrac2n + O(n^{-2})\bigr)$ ---
so for $n$ in the hundreds the Poisson model is already
accurate to the third digit. This is the practical content
of the law of rare events: the modeller never knows $n$ and
$p$ separately (how many micro-opportunities for a typo does
a page hold?), but only their product $\lambda$, and the
limit law mercifully depends on nothing else.
\end{example}

\section{Branching processes}

Consider a population starting from one ancestor; each individual,
independently, has a random number of children with law $(p_k)_{k
\in \N}$ and generating function $G$ (the \emph{offspring
distribution}). Let $Z_n$ be the size of generation $n$ ($Z_0 =
1$), and let $m = G'(1) = \E(Z_1)$ be the mean offspring number.

\begin{proposition}\label{prop:b2:genfun:branching}
The generating function of $Z_n$ is the $n$-th iterate
$G_{Z_n} = G \circ G \circ \dots \circ G$ ($n$ times), and the
extinction probabilities $q_n = \P(Z_n = 0)$ satisfy
\[
q_0 = 0, \qquad q_{n+1} = G(q_n),
\]
and increase to the probability $q$ of eventual extinction, which
is a fixed point of $G$.
\end{proposition}

\begin{proof}
Generation $n + 1$ is the random sum of the offspring of the $Z_n$
members of generation $n$, with counts independent of each other
and of $Z_n$: \cref{thm:b2:genfun:compound} gives $G_{Z_{n+1}} =
G_{Z_n} \circ G$, and induction from $G_{Z_0}(t) = t$ yields the
$n$-fold iterate --- which, by associativity of composition, can
equally be read as $G_{Z_{n+1}} = G \circ G_{Z_n}$. Evaluating
this second form at $0$: $q_{n+1} = G_{Z_{n+1}}(0) =
G\bigl(G_{Z_n}(0)\bigr) = G(q_n)$. The events $\{Z_n = 0\}$ increase (extinct populations stay
extinct), so $q_n \uparrow q = \P\bigl(\bigcup_n\{Z_n = 0\}\bigr)$
by monotone continuity (\cref{thm:b2:proba:continuity}), and
continuity of $G$ on $[0, 1]$ turns $q_{n+1} = G(q_n)$ into $q =
G(q)$ at the limit.
\end{proof}

\begin{example}[Watching extinction converge]
\label{ex:b2:genfun:cobwebnumerics}
For the offspring law $(p_0, p_1, p_2) = (\tfrac14,
\tfrac14, \tfrac12)$ of \cref{ex:b2:genfun:branchingexample},
$G(t) = \tfrac14 + \tfrac14t + \tfrac12t^2$ and the iteration
$q_{n+1} = G(q_n)$ gives
\[
q_1 = 0.25, \quad q_2 = 0.34375, \quad q_3 \approx 0.39502,
\quad q_4 \approx 0.42678, \quad q_5 \approx 0.44776,
\]
climbing toward the extinction probability $q = \tfrac12$.
The gaps $q - q_n$ are $0.25$, $0.156$, $0.105$, $0.073$,
$0.052$: each is roughly $\tfrac34$ of the previous, and
indeed the mean value theorem gives $q - q_{n+1} = G'(c_n)(q
- q_n)$ with $G'(q) = \tfrac14 + q = \tfrac34$. Two morals: a
family line still alive at generation $n$ has, built into
the same computation, probability $q - q_n$ of being doomed
later; and the convergence rate of the staircase in the
figure below is the derivative at the fixed point --- the
weekend problem turns both observations into theorems.
\end{example}

\begin{theorem}[Extinction criterion]\label{thm:b2:genfun:extinction}
Assume $p_1 \neq 1$. The extinction probability $q$ is the
\emph{smallest} fixed point of $G$ in $\intcc{0}{1}$, and:
\begin{itemize}
\item if $m \leq 1$ (subcritical or critical), $q = 1$: extinction
is certain;
\item if $m > 1$ (supercritical), $q < 1$: the population survives
forever with positive probability $1 - q$.
\end{itemize}
\end{theorem}

\begin{proof}
$G$ is convex on $\intcc{0}{1}$ (power series with nonnegative
coefficients: $G'' \geq 0$), nondecreasing, with $G(1) = 1$.

\emph{Smallest fixed point:} let $r \in \intcc{0}{1}$ be any fixed
point. Then $q_0 = 0 \leq r$, and inductively $q_{n+1} = G(q_n)
\leq G(r) = r$ (monotonicity): so $q = \lim q_n \leq r$.

\emph{Case $m \leq 1$:} suppose $r < 1$ is a fixed point. By the
mean value theorem on $[r, 1]$, there is $c \in \intoo{r}{1}$ with
$G'(c) = \frac{G(1) - G(r)}{1 - r} = \frac{1 - r}{1 - r} = 1$. But
$G'$ is nondecreasing (convexity) with $\lim_{t\to1^-}G'(t) = m
\leq 1$, so $G' \leq 1$ on $\intoo{0}{1}$; the equality $G'(c) =
1$ then forces $G'$ to be constant equal to $1$ on $\intco{c}{1}$,
hence $G'' = \sum n(n-1)p_nt^{n-2} \equiv 0$ there. A power series
with nonnegative coefficients vanishing on an interval has all
these coefficients zero: $p_n = 0$ for $n \geq 2$, so $G(t) = p_0
+ p_1t$ and $1 = G'(c) = p_1$ --- contradicting the hypothesis
$p_1 \neq 1$. So $1$ is the only fixed point: $q = 1$.

\emph{Case $m > 1$:} near $1$, $G(t) - t$ has derivative $G'(t) -
1 \to m - 1 > 0$ as $t \to 1^-$, so $G(t) - t < G(1) - 1 = 0$ on
some interval $\intoo{1 - \delta}{1}$: the continuous function
$G(t) - t$ is $\geq 0$ at $t = 0$ ($G(0) = p_0 \geq 0$) and $< 0$
just below $1$, so it vanishes at some $r < 1$ (intermediate value
theorem). The smallest fixed point is then $q \leq r < 1$.
\end{proof}

\begin{figure}[ht]
\centering
% Subcritical: G(t) = 0.5 + 0.1 t + 0.4 t^2 => m = 0.9 <= 1
\begin{tikzpicture}[scale=2.9]
  \draw[->] (0,0) -- (1.15,0) node[below] {$t$};
  \draw[->] (0,0) -- (0,1.15);
  \draw[dashed] (0,0) -- (1,1);
  \draw[thick, ocBlue, domain=0:1, smooth, samples=40]
    plot (\x, {0.5 + 0.1*\x + 0.4*\x*\x});
  \fill (1,1) circle (0.02);
  \node[below, font=\small] at (0.55,-0.06)
    {$m \leq 1$: $q = 1$};
  % cobweb from 0
  \draw[ocRed, ->] (0,0) -- (0,0.5);
  \draw[ocRed, ->] (0,0.5) -- (0.5,0.5);
  \draw[ocRed, ->] (0.5,0.5) -- (0.5,0.65);
  \draw[ocRed, ->] (0.5,0.65) -- (0.65,0.65);
  \draw[ocRed, ->] (0.65,0.65) -- (0.65,0.734);
  \draw[ocRed, ->] (0.65,0.734) -- (0.734,0.734);
  \node[font=\small, ocBlue] at (0.32,0.75) {$G$};
\end{tikzpicture}
\qquad
% Supercritical: G(t) = 0.25 + 0.15 t + 0.6 t^2 => m = 1.35 > 1
\begin{tikzpicture}[scale=2.9]
  \draw[->] (0,0) -- (1.15,0) node[below] {$t$};
  \draw[->] (0,0) -- (0,1.15);
  \draw[dashed] (0,0) -- (1,1);
  \draw[thick, ocBlue, domain=0:1, smooth, samples=40]
    plot (\x, {0.25 + 0.15*\x + 0.6*\x*\x});
  \fill (1,1) circle (0.02);
  % fixed point q: 0.6q^2 - 0.85q + 0.25 = 0 -> q = 5/12 ~ 0.4167
  \fill[ocRed] (0.4167,0.4167) circle (0.02)
    node[below right, font=\small] {$q$};
  \node[below, font=\small] at (0.55,-0.06)
    {$m > 1$: $q < 1$};
  % cobweb from 0
  \draw[ocRed, ->] (0,0) -- (0,0.25);
  \draw[ocRed, ->] (0,0.25) -- (0.25,0.25);
  \draw[ocRed, ->] (0.25,0.25) -- (0.25,0.325);
  \draw[ocRed, ->] (0.25,0.325) -- (0.325,0.325);
  \draw[ocRed, ->] (0.325,0.325) -- (0.325,0.3623);
  \draw[ocRed, ->] (0.325,0.3623) -- (0.3623,0.3623);
  \node[font=\small, ocBlue] at (0.28,0.72) {$G$};
\end{tikzpicture}
\caption{Extinction probabilities as a fixed-point iteration
$q_{n+1} = G(q_n)$ starting at $q_0 = 0$ (red staircase). Left: a
subcritical offspring law --- the convex curve stays above the
diagonal, the iteration climbs to the unique fixed point $1$.
Right: a supercritical law --- the curve crosses the diagonal at
$q < 1$, where the iteration stops: survival has probability $1 -
q > 0$.}
\label{fig:b2:genfun:branching}
\end{figure}

\begin{remark}[How to read the cobweb]
In the figure, a vertical move applies $G$ (from $(q_n, q_n)$
up to $(q_n, G(q_n))$), a horizontal move to the diagonal
converts output into input: the staircase \emph{is} the
recursion $q_{n+1} = G(q_n)$. Convexity of $G$ and $G(1) = 1$
leave only two geometries. Either the curve stays above the
diagonal on $\intco01$ (mean $m \leq 1$): the staircase has
nowhere to stop before $1$. Or the curve crosses at some $q <
1$ ($m > 1$): the staircase is trapped below the crossing
and converges to it, at the geometric rate $G'(q) < 1$
quantified in \cref{ex:b2:genfun:cobwebnumerics}. All the
analysis of the extinction theorem is visible in this one
picture --- which is why it is worth drawing before
computing.
\end{remark}

\begin{example}\label{ex:b2:genfun:branchingexample}
Offspring law: no child, one child, two children with
probabilities $\frac14, \frac14, \frac12$. Then $m = \frac14 + 1 =
\frac54 > 1$ and $G(t) = \frac14 + \frac14 t + \frac12 t^2$. Fixed
points: $\frac12 t^2 - \frac34 t + \frac14 = 0$, i.e.\ $2t^2 - 3t
+ 1 = (2t - 1)(t - 1) = 0$: $q = \frac12$. The family line dies
out with probability $\frac12$ --- and with probability $\frac12$
it lives forever.
\end{example}

\begin{remark}[Perspectives within this volume]
The chapter is the book's crossroads, and each ingredient
arrived from a named place: the series algebra from
\cref{ch:b2:series} and \cref{ch:b2:powerseries}, the
probability from \cref{ch:b2:proba} (monotone continuity
proves $q_n \uparrow q$) and \cref{ch:b2:randomvar} ($G_X =
\E(t^X)$ is an expectation, multiplicativity is the product
theorem), the convexity from \cref{ch:b2:realfun} through
\cref{ch:b2:affine}. Even the heavy-tail pathologies connect:
the St Petersburg variable of the previous chapter has
$G(t) = \sum_k2^{-k}t^{2^k}$, a perfectly convergent series
on $\intcc01$ whose derivative at $1^-$ diverges ---
infinite mean, visible at a glance. One object, every tool
of the year: a fitting last chapter.
\end{remark}

\begin{remark}[Common pitfalls]
(i) Generating functions apply to $\N$-valued variables only:
for signed or non-integer variables the object $\E(t^X)$
loses its power-series structure (Year 3 replaces it with
transforms adapted to $\R$). (ii) The first sanity check of
any computed $G$ is $G(1) = 1$; the second is that the
coefficients are nonnegative --- a negative coefficient
means an algebra slip, not a new law. (iii) In random sums,
the composition order matters: $G_S = G_N \circ G_X$, the
\emph{outer} function counting the terms; composing the
other way is meaningless ($G_X \circ G_N$ would count items
of items). (iv) Multiplicativity needs independence and
distinct sources of randomness: $G_{2X}(t) = G_X(t^2)$, not
$G_X(t)^2$. (v) Differentiating at $1$ is a boundary
operation: when the radius is exactly $1$, as in
\cref{ex:b2:genfun:heavytail}, $G'(1^-)$ may be infinite,
and the moments theorem's monotone-limit formulation is not
a pedantic nicety but the honest statement.
\end{remark}

\section*{Closing the volume}

The generating function is a fitting final object for this book:
it is simultaneously a power series (\cref{ch:b2:powerseries}), a
tool of summable families (\cref{ch:b2:series}), an expectation
(\cref{ch:b2:randomvar}), a convex function whose geometry decides
extinction (\cref{ch:b2:realfun}), and a fixed-point iteration
(\cref{ch:b2:metric}). The mathematics of Year 2 is one subject.
The Year 3 volume will open the doors deliberately left closed
here: Lebesgue integration (discharging the dominated convergence
theorem of \cref{ch:b2:integration}), measure-theoretic
probability on uncountable spaces, and the inverse function
theorem's full proof (\cref{ch:b2:diffcalc}) in the setting of
differential geometry.

\section{Exercises}

\begin{exercise}[$\star$]\label{exo:b2:genfun:1}
Compute the generating function of the uniform law on $\{1, 2,
\dots, 6\}$ (a fair die). Show that the sum of two fair dice
\emph{cannot} be uniform on $\{2, \dots, 12\}$: factor $G_{X+Y}$
and count roots. \emph{(A uniform sum would force $G_X(t)G_Y(t) =
\frac{t^2}{11}\sum_{k=0}^{10}t^k$, whose nonzero roots are the
$11$-th roots of unity other than $1$ --- none of them real ---
while $G_X/t$ and $G_Y/t$ are real polynomials of degree $5$, each
owning at least one real root.)}
\end{exercise}

\begin{exercise}[$\star$]\label{exo:b2:genfun:2}
Using generating functions, recover $\E$ and $V$ for the binomial
and geometric laws (\cref{thm:b2:genfun:moments}).
\end{exercise}

\begin{exercise}[$\star$]\label{exo:b2:genfun:3}
Two loaded dice: is it possible to load two dice (independently,
identically or not) so that their sum is uniform on $\{2, \dots,
12\}$? \emph{(Same factorization obstruction as in
\cref{exo:b2:genfun:1}: the answer is no even with different
loadings, because each factor $G_X(t)/t$ has odd degree $5$,
hence a real root, while the target has none.)}
\end{exercise}

\begin{exercise}[$\star\star$]\label{exo:b2:genfun:4}
Let $X_1, X_2, \dots$ be independent Bernoulli $\mathcal{B}(p)$
and $N \sim \mathcal{P}(\lambda)$ independent of them. Show,
via \cref{thm:b2:genfun:compound}, that $S = X_1 + \dots + X_N
\sim \mathcal{P}(\lambda p)$: a Poisson number of items, each kept
with probability $p$, leaves a Poisson number --- \emph{thinning}.
Compute also the law of the discarded count and admire: it is
$\mathcal{P}(\lambda(1-p))$, and one can show it is independent
of $S$.
\end{exercise}

\begin{exercise}[$\star\star$]\label{exo:b2:genfun:5}
(Negative binomial) Let $T_r$ be the number of tosses to obtain
$r$ heads (head probability $p$). Write $T_r$ as a sum of $r$
independent geometric variables, deduce
\[
G_{T_r}(t) = \Bigl(\frac{pt}{1 - (1-p)t}\Bigr)^{r},
\qquad
\E(T_r) = \frac rp,
\qquad
V(T_r) = \frac{r(1-p)}{p^2},
\]
and expand $G_{T_r}$ to find $\P(T_r = n) = \binom{n-1}{r-1}
p^r(1-p)^{n-r}$.
\end{exercise}

\begin{exercise}[$\star\star$]\label{exo:b2:genfun:6}
For the offspring law $p_0 = \frac18$, $p_1 = \frac38$, $p_2 =
\frac38$, $p_3 = \frac18$: compute $m$, decide
supercriticality, and compute the extinction probability $q$
exactly. \emph{(Factor out the root $t = 1$ of $G(t) - t$.)}
\end{exercise}

\begin{exercise}[$\star\star\star$]\label{exo:b2:genfun:7}
(Total progeny) In a subcritical branching process ($m < 1$), let
$Y = \sum_{n\geq0} Z_n$ be the total number of individuals ever
born. Show $\E(Y) = \sum_n m^n = \frac{1}{1 - m}$ (justify the
exchange of summations), and prove that the generating function
$H = G_Y$ satisfies the functional equation $H(t) = t\,G(H(t))$.
\emph{(The ancestor, plus the total progenies of each of its
children, which are independent copies of $Y$.)}
\end{exercise}

\begin{exercise}[$\star\star\star$]\label{exo:b2:genfun:8}
Let $X$ have generating function $G$ with radius of convergence $>
1$. Prove the \emph{exponential tail bound}: there are $C > 0$ and
$\rho \in \intoo{0}{1}$ with $\P(X \geq n) \leq C\rho^n$.
\emph{(Markov applied to $t^X$ for a fixed $t > 1$ inside the
disk.)} Conversely, show that if $\P(X \geq n) \leq C\rho^n$ with
$\rho < 1$, the radius of $G$ is $\geq 1/\rho > 1$.
\end{exercise}

\begin{exercise}[$\star\star\star$]\label{exo:b2:genfun:9}
(Continuity theorem, elementary case) Let $X, X_1, X_2, \dots$ be
$\N$-valued with $G_{X_n}(t) \to G_X(t)$ for every $t \in
\intco{0}{1}$. Show that $\P(X_n = k) \to \P(X = k)$ for every
$k$. \emph{(Induction on $k$: for $k = 0$ take $t \to 0$ ---
carefully: fix $t$ small, use $\abs{\P(X_n = 0) - G_{X_n}(t)}
\leq \frac{t}{1-t}$, valid since the tail $\sum_{j \geq 1}p_jt^j
\leq \frac{t}{1 - t}$; then diagonalize. For the induction step,
consider $\frac{G(t) - \P(X = 0)}{t}$, the generating function of
a shifted law.)}
\end{exercise}

\begin{exercise}[$\star$]\label{exo:b2:genfun:10}
(Parity trick) Show that for an $\N$-valued variable $X$,
\[
\P(X \text{ even}) = \frac{1 + G_X(-1)}{2} ,
\]
and compute this probability for $X \sim \mathcal P(\lambda)$
and $X \sim \mathcal B(n, p)$. What does $G_X(-1) \to 0$ mean
probabilistically?
\end{exercise}

\begin{exercise}[$\star\star$]\label{exo:b2:genfun:11}
(Sicherman dice) Verify the factorization of the fair-die
generating function
\[
\frac{t + t^2 + \dots + t^6}{6}
= \frac{t\,(1 + t)(1 + t + t^2)(1 - t + t^2)}{6},
\]
and show that the two dice with faces $\{1, 2, 2, 3, 3, 4\}$
and $\{1, 3, 4, 5, 6, 8\}$ have generating functions
$\frac{t(1+t)(1+t+t^2)}6$ and
$\frac{t(1+t)(1+t+t^2)(1-t+t^2)^2}6$, whose product is that of
two standard dice: these exotic dice produce every total $2,
\dots, 12$ with exactly the standard probabilities.
\end{exercise}

\begin{exercise}[$\star\star\star$]\label{exo:b2:genfun:12}
(Waiting for two heads in a row) A coin with head probability
$p$ is tossed until two consecutive heads appear; let $T$ be
the number of tosses (the game of \cref{exo:b2:proba:6}).
Conditioning on the first tosses, derive a linear system for
the generating functions from the states ``no current head''
and ``one current head'', and conclude
\[
G_T(t) = \frac{p^2t^2}{1 - qt - pqt^2}
\qquad (q = 1 - p);
\]
check $G_T(1) = 1$ and $\E(T) = \dfrac{1 + p}{p^2}$ ($= 6$
for a fair coin).
\end{exercise}

\section{Problem: the Galton--Watson process, solved}

\begin{problem}[{Weekend problem --- growth rates, exact
solutions, total progeny, and Kolmogorov's critical
estimate}]\label{pb:b2:genfun:1}
The extinction criterion (\cref{thm:b2:genfun:extinction})
splits branching processes\index{branching process} into
subcritical, critical and supercritical --- but it says
nothing about \emph{rates}: how fast a doomed line dies, how
big a surviving one grows. This problem computes them. We keep
the chapter's notation: offspring law $(p_k)$ with pgf $G$,
mean $m = G'(1)$, generation sizes $Z_n$ ($Z_0 = 1$), iterates
$G_n = G_{Z_n}$, extinction probabilities $q_n = \P(Z_n = 0)
\uparrow q$; we always assume $p_1 \neq 1$ and, where second
moments appear, $G''(1) < \infty$, and we write $\sigma^2 =
V(Z_1)$.

\medskip
\noindent\textbf{Part I --- Moments of the generations.}
\begin{enumerate}
  \item Show $\E(Z_n) = m^n$ \emph{(chain rule on $G_n = G
        \circ G_{n-1}$ at $1^-$, using $G_{n-1}(1) = 1$ and
        \cref{thm:b2:genfun:moments})}.
  \item Establish the recursion $G_n''(1) =
        G''(1)\,m^{2(n-1)} + m\,G_{n-1}''(1)$ and solve it:
        $G_n''(1) = G''(1)\,m^{n-1}\dfrac{m^n - 1}{m - 1}$
        for $m \neq 1$, and $G_n''(1) = n\,G''(1)$ for $m =
        1$.
  \item Deduce
        \[
        V(Z_n) = \sigma^2m^{n-1}\,\frac{m^n - 1}{m - 1}
        \quad (m \neq 1),
        \qquad
        V(Z_n) = n\,\sigma^2 \quad (m = 1).
        \]
  \item (Subcritical rate, upper bound) For $m < 1$, show
        $\P(Z_n > 0) \leq m^n$ \emph{(Markov on the
        integer-valued $Z_n$)}: extinction is certain with a
        geometric rate --- a quantitative refinement of the
        chapter's criterion.
  \item (Subcritical rate, lower bound) Using
        Cauchy--Schwarz on $Z_n\mathbf 1_{Z_n > 0}$, show
        \[
        \P(Z_n > 0) \geq \frac{\E(Z_n)^2}{\E(Z_n^2)}
        \geq c\,m^{n}
        \quad\text{with}\quad
        c = \Bigl(\frac{\sigma^2}{m(1-m)} + 1\Bigr)^{-1} :
        \]
        the geometric rate $m^n$ is exact up to constants.
\end{enumerate}

\medskip
\noindent\textbf{Part II --- The geometric family, solved
exactly.} Let the offspring law be geometric on $\N$: $p_k =
qp^k$ ($k \geq 0$), with $0 < p < 1$, $q = 1 - p$.
\begin{enumerate}[resume]
  \item Compute $G(t) = \dfrac{q}{1 - pt}$ and $m = \dfrac
        pq$; locate the three regimes in terms of $p$.
  \item Solve $G(t) = t$: show the fixed points are $1$ and
        $q/p = 1/m$, and recover the extinction probability
        $q_{\mathrm{ext}} = \min(1, 1/m)$.
  \item Prove by induction the closed forms
        \[
        q_n = \frac{m^n - 1}{m^{n+1} - 1} \quad (m \neq 1),
        \qquad
        q_n = \frac{n}{n+1} \quad (m = 1).
        \]
  \item Deduce the exact rates: $1 - q_n \sim (1 - m)\,m^n$
        in the subcritical case, and $q_{\mathrm{ext}} - q_n
        \sim \dfrac{m - 1}{m^{2}}\cdot m^{-n}$ in the
        supercritical case; check that the supercritical
        contraction ratio is $G'(q_{\mathrm{ext}}) = 1/m$.
  \item Critical case ($p = \tfrac12$): compute $\sigma^2 =
        2$ and note $1 - q_n = \frac1{n+1}$: survival decays
        like $\frac1n$ --- neither geometric nor summable.
  \item Still critical: prove by induction the full iterate
        \[
        G_n(t) = \frac{n - (n-1)t}{n + 1 - nt},
        \]
        and deduce that conditioned on survival, $Z_n$ is
        geometric on $\N^*$ with parameter $\frac1{n+1}$:
        \[
        \P(Z_n = k \mid Z_n > 0) = \frac1{n+1}
        \Bigl(\frac{n}{n+1}\Bigr)^{k-1},
        \qquad
        \E(Z_n \mid Z_n > 0) = n + 1 .
        \]
        The average line dies, but the surviving lines have
        size of order $n$.
\end{enumerate}

\medskip
\noindent\textbf{Part III --- Total progeny.} Let $Y =
\sum_{n\geq0}Z_n \in \N^* \cup \{\infty\}$ be the total number
of individuals ever born, and $H(t) =
\sum_{k\geq1}\P(Y = k)t^k$.
\begin{enumerate}[resume]
  \item Justify $\P(Y < \infty) = q_{\mathrm{ext}}$, and
        recall from \cref{exo:b2:genfun:7} the functional
        equation $H(t) = t\,G(H(t))$ (whose derivation did
        not use $m < 1$).
  \item (Binary branching) For $p_0 = p_2 = \frac12$
        (critical), solve the functional equation:
        \[
        H(t) = \frac{1 - \sqrt{1 - t^2}}{t},
        \]
        and expand with \cref{ex:b2:powerseries:catalan} to
        get
        \[
        \P(Y = 2k + 1) = \frac{C_k}{2^{2k+1}},
        \qquad C_k = \frac1{k+1}\binom{2k}k ;
        \]
        check the values $\P(Y = 1) = \frac12$ and $\P(Y = 3)
        = \frac18$ by direct counting.
  \item Differentiating the functional equation at $1^-$,
        show that $\E(Y) = \frac{1}{1-m}$ for $m < 1$, while
        criticality forces $\E(Y) = \infty$: the critical
        total progeny is finite almost surely with infinite
        mean.
  \item With the central binomial asymptotics
        (\cref{ex:b2:comparison:centralbinomial}), show
        \[
        \P(Y = 2k+1) \sim \frac{1}{2\sqrt\pi\,k^{3/2}},
        \]
        a heavy $k^{-3/2}$ tail, and deduce $\P(Y > n)
        \asymp n^{-1/2}$ (upper and lower bounds of this
        order suffice).
  \item Compare with the fair random walk (the weekend
        problem of \cref{ch:b2:proba}): certain but
        infinite-mean return times there, certain but
        infinite-mean total progeny here, both with
        $n^{-3/2}$ local laws. One paragraph on why
        criticality produces this signature.
\end{enumerate}

\medskip
\noindent\textbf{Part IV --- Kolmogorov's estimate at
criticality.} Assume $m = 1$, $0 < \sigma^2 = G''(1) <
\infty$.
\begin{enumerate}[resume]
  \item Show that $G''$ extends continuously to $\intcc01$
        \emph{(nonnegative increasing with finite limit)} and
        deduce the Taylor expansion at $1$:
        \[
        G(t) = t + b\,(1-t)^2 + o\bigl((1-t)^2\bigr),
        \qquad b = \frac{G''(1)}2 = \frac{\sigma^2}2 .
        \]
  \item For $t \in \intco01$, set $h(t) = \dfrac1{1 - G(t)} -
        \dfrac1{1 - t}$. Show
        \[
        h(t) = \frac{G(t) - t}{(1 - G(t))(1 - t)}
        \xrightarrow[t\to1^-]{} b .
        \]
  \item Telescope along the iteration $q_{j+1} = G(q_j)$:
        \[
        \frac1{1 - q_n} = 1 + \sum_{j=0}^{n-1}h(q_j),
        \]
        and conclude with a Cesaro argument that
        \[
        \P(Z_n > 0) = 1 - q_n \sim \frac{2}{\sigma^2\,n}
        \]
        --- \emph{Kolmogorov's estimate}: every critical
        branching process dies at the universal rate $1/n$,
        with only the constant remembering the offspring law.
  \item Check the estimate against the critical geometric
        case of question 10.
  \item Deduce $\E(Z_n \mid Z_n > 0) = \dfrac{1}{1 - q_n}
        \sim \dfrac{\sigma^2 n}{2}$ \emph{(note $\E(Z_n
        \mathbf 1_{Z_n>0}) = \E(Z_n) = 1$)}, and check it
        against question 11: conditioned on survival, the
        population grows \emph{linearly} --- the critical
        tightrope between death and explosion.
\end{enumerate}

\medskip
\noindent\textbf{Part V --- Applications and synthesis.}
\begin{enumerate}[resume]
  \item (Epidemics, chain reactions) For a Poisson offspring
        law $\mathcal P(\lambda)$ --- each case infects
        $\mathcal P(\lambda)$ new cases --- write the
        extinction equation $q = \eu^{\lambda(q-1)}$ and
        solve it numerically for $\lambda = 1.5$ ($q \approx
        0.417$) and $\lambda = 2$ ($q \approx 0.203$):
        starting from one case, a major outbreak is
        \emph{not} certain even when $\lambda > 1$. Explain
        why the iteration $q_{n+1} = \eu^{\lambda(q_n - 1)}$
        from $q_0 = 0$ converges to the right root.
  \item Starting from $k$ ancestors instead of one, show that
        the extinction probability is $q^k$. Application:
        with $\lambda = 1.5$, how many initial cases make an
        outbreak at least $99\%$ likely?
  \item (Conditioning a supercritical process on extinction)
        For $m > 1$ with extinction probability $q \in
        \intoo01$: first prove by convexity that $G'(q) < 1$
        at the smallest fixed point, and deduce
        $q_{\mathrm{ext}} - q_n = O\bigl(G'(q)^n\bigr)$
        (geometric convergence, as question 9 exemplified).
        Then show that $\widehat G(t) = G(qt)/q$ is the pgf
        of a bona fide offspring law, with mean $\widehat m =
        G'(q) < 1$: a subcritical companion process. Verify
        on the geometric family: conditioning the
        supercritical $(p, q)$-process on extinction swaps
        $p$ and $q$. (The full statement --- the conditioned
        process \emph{is} the companion process --- is proved
        in the Year 3 volume; here you have verified its
        generating-function shadow.)
  \item Synthesis: draw up the trichotomy table --- for $m <
        1$, $m = 1$, $m > 1$: value of $q$; rate of $\P(Z_n >
        0)$ or of $q - q_n$; $\E(Y)$; size of a surviving
        generation. State in one sentence per tool how
        composition of pgfs, convexity, Taylor at $1^-$, and
        Cesaro averaging carried the whole problem, and what
        the Year 3 volume adds (the martingale $Z_n/m^n$ and
        Yaglom's exponential limit law).
\end{enumerate}
\end{problem}
